\sscp{Clause-final negation}{02-03-01-02}
The position of standard negation within a clause is another
typological feature for which contact-induced change provides
the best explanation for certain patterns found in the region.
But like the preposed possessors,
those typological patterns which are the result of contact should
not be used as a basis for assigning linguistic affiliation.
The dominant western Austronesian pattern for standard negation is pre-verbal.
Clause-final negation in the region is associated with verb-final SOV Papuan languages.
\citep{Reesink2002b,KlamerReesinkVanStaden2008,Florey2010,Dryer2013e,ReesinkDunn2018}.
 Yet in our region we find branches of several subgroups have
within them languages with negators that are pre-verbal,
clause-final, or have bipartite (split, embracing) negation.
For example, the \tsc{Flores-Lembata} node of \tsc{\ili{Bima}-Lembata}
\citep{Fricke2017}, \tsc{\ili{Ambon}-Seram} \citep{Florey2010}, and the \tsc{Rote-Meto}
node of \tsc{Timor-Babar} (illustrated below) each have all three
patterns of standard negation within them.

Clause-final negation is common in several Austronesian subgroups
in this region, and is also found in SHWNG languages.
\citet{GrimesCIP} points out that in this region bipartite
negation provides a kind of missing link -- a mechanism to transition
from pre-verbal to clause-final negation.

\ili{Malay}, in example \qf{02-22}, illustrates the
common pattern for Austronesian
languages west of the Blust line,
placing standard negation before the verb.

\begin{exe}
\setcounter{xnumi}{21}
	\ex{\ili{Malay}; western pattern; preverbal}\vspace{-1mm}
	\sn{\gll	saya \tbf{tidaʔ} tahu\\
						\tsc{1sg} \tsc{neg} know\\
			\glt	`I don't know.'
						}\label{ex:02-22}
\end{exe}

Examples \qf{02-23}--\qf{02-23b} illustrate preverbal,
bipartite, and clause-final negation from \tsc{Timor-Babar}
languages.

\begin{exe}
	\ex{\ili{Tetun} \ili{Foho} [tet]; preverbal\label{ex:02-23}}\jambox*{\citep[751]{UBB2013}}\vspace{-1mm}
	\sn{\gll	haʔu \tbf{la} k-atene\\
						\tsc{1sg} \tsc{neg} \tsc{1sg}-know\\
			\glt	`I don't know.'}
	\ex{\ili{Dela} [row]; bipartite\label{ex:02-24}}\jambox*{\citep[945]{UBB2020a}}\vspace{-1mm}
	\sn{\gll	au \tbf{nda} ɓuɓuluʔ \tbf{sa}\\
						\tsc{1sg} \tsc{neg1} know \tsc{neg2}\\
			\glt	`I don't know.'}
\end{exe}

\newpage
\begin{exe}
	\ex{\ili{Tetun} \ili{Fehan} [tet]; clause final\label{ex:02-23b}}\jambox*{(Grimes' fieldnotes)}\vspace{-1mm}
	\sn{\gll	ha k-atene \tbf{haʔi}\\
						\tsc{1sg} \tsc{1sg}-know \tsc{neg} \\
			\glt	`I don't know.'}
	%\ex{\ili{Helong} [heg]; clause final\label{ex:02-25}}\jambox*{\citep[855]{UBB2012a}}\vspace{-1mm}
	%\sn{\gll	auk tanan \tbf{lo}\\
						%\tsc{1sg} know \tsc{neg}\\
			%\glt	`I don't know.'}
\end{exe}

While \ili{Tetun}, \ili{Dela}, and \ili{Helong} are separated from each other by
other islands or languages,
it should be noted that within the \tsc{Meto} cluster of closely
related dialects and languages itself (which are also \tsc{Timor-Babar} languages)
we find similar variation.
For example, \ili{Kotos} \ili{Amarasi} has bipartite negation \it{ka {\ldots} fa}.
Neighbouring Roi{\Q}is \ili{Amarasi} has clause final negation \it{maeʔ} \citep[412]{Edwards2020}.
Recorded texts from younger speakers of Amfo{\Q}an often have only
the pre-verbal \it{ka} with no corresponding \it{fa}.
A similar trend of dropping the second part of bipartite negation
is found among younger speakers of Rote-cluster languages \ili{Tii} and \ili{Lole}.
Both parts tend to be preserved by some older speakers.
This trend is assumed to be influenced from Indonesian with its
single preverbal standard negation.

Examples \qf{02-26}--\qf{02-28} illustrate preverbal,
bipartite, and clause-final negation from three \tsc{\ili{Ambon}-Seram} languages.

\begin{exe}
	\ex{Allang [alo]; preverbal\label{ex:02-26}}\jambox*{\citep[233]{Florey2010}}\vspace{-1mm}
	\sn{\gll	nalisa ite \tbf{ta} supu luma-nu\\
						tomorrow \tsc{1pl.excl} \tsc{neg} meet \tsc{recip-noun.class.marker}\\
			\glt	`We won't meet each other tomorrow.'}
	\ex{\ili{Haruku} [hrk]; bipartite\label{ex:02-27}}\jambox*{\citep[233]{Florey2010}}\vspace{-1mm}
	\sn{\gll	au \tbf{taʔ} ane \tbf{sa} hahu\\
						\tsc{1sg} \tsc{neg1} eat \tsc{neg2} pig\\
			\glt	`I don't eat pork.'}
	\ex{\ili{Yalahatan} [jal]; clause final\label{ex:02-28}}\jambox*{\citep{BeckleyND}}\vspace{-1mm}
	\sn{\gll	yaʔu nia-i mo mosa-i \tbf{tam}\\
						\tsc{1sg} look.for-\tsc{3sg.u} but find-\tsc{3sg.u} \tsc{neg}\\
			\glt	`I looked for her but didn't find her.'}
\end{exe}

Examples \qf{02-23}--\qf{02-28} above show that these very different structural
patterns can co-exist in neighbouring languages within the same
branch of the same subgroup,
and even within the same language-dialect cluster.
The Jespersen Cycle of grammaticalisation shows that languages
can change diachronically from pre-predicate negation to bipartite
negation to post-predicate negation
\citep{Jespersen1917,Vossen2016,Fricke2017}.
But the Jespersen Cycle does not explain why languages west of
the Blust line tend to be limited to pre-verbal negation,
whereas languages within our zone of contact that also show traces
of a Papuan substrate in other features,
are the ones that show all three patterns of negation.
Contact with Papuan languages that have bipartite
and clause-final negation provides the trigger or pressure that
allows the Jespersen Cycle to operate.
Why did the shift in negation occur so regularly within this
region, but not west of it? Why does clause-final negation occur in so
many SVO languages \ili{east} of the Blust line,
but not west of it?

To assume that clause-final negation in so many SVO Austronesian
languages in the region developed from the pre-verbal negation
associated cross-linguistically with SVO typology
and pervasive west of the Blust line,
as an innovation in a vacuum is not a reasonable explanation.
It is far more reasonable to attribute the clause-final negation
in Austronesian languages to contact with SOV Papuan languages
that already have that pattern associated cross-linguistically with SOV typology.
And its presence is seen as trace evidence of intense historical
contact, even where the nearest Papuan languages spoken today are hundreds
of kilometres away, such as in the Austronesian languages of \ili{Dhao},
\ili{Helong}, \ili{Buru}, and many languages on the island of Seram.


